
Deep neural networks achieve high average accuracy on benchmark datasets but frequently fail on minority subgroups when spurious correlations are present. On the Waterbirds dataset, which contains bird images with spurious background correlations, standard empirical risk minimization (ERM) training achieves 89.2\% average accuracy but only 72.3\% accuracy on the worst-performing group. In contrast, Group Distributionally Robust Optimization (Group-DRO), which explicitly minimizes worst-group loss during training, achieves 88.7\% worst-group accuracy. This performance gap indicates that gradient-based optimization preferentially exploits spurious features rather than learning core discriminative features.

The causes of this optimization behavior are not well understood. Existing work has focused on developing training methods that improve worst-group performance, including Group-DRO, which requires group labels during training, and Sharpness-Aware Minimization (SAM), which seeks flat minima but does not distinguish spurious from core features. These approaches provide solutions without explaining why standard optimization fails.

This study investigates whether loss landscape geometry differentiates models trained with standard versus robust methods. Specifically, we examine whether Hessian eigenvalue structure, analyzed through Marchenko-Pastur random matrix theory, provides a geometric signature of spurious correlation exploitation. We measure the alignment between minority-group gradients and sharp curvature directions defined by outlier eigenvalues exceeding the Marchenko-Pastur bulk edge.

Our experiments on Waterbirds demonstrate three findings. First, ERM solutions exhibit substantially higher alignment between minority gradients and sharp curvature subspaces compared to Group-DRO solutions (72.3\% vs. 31.6\%, Cohen's d = 1.87). Second, ERM produces 53\% more outlier eigenvalues than Group-DRO, indicating concentrated sharp curvature. Third, stochastic gradient descent trajectories show directional bias toward flat directions (0.15, p = 0.023). However, two proposed mechanism components could not be validated: minority gradient alignment to real Hessian eigenvectors was not computed, and geometric alignment did not correlate with worst-group accuracy along interpolation paths.

These results are limited to a single dataset (Waterbirds) with one random seed for the primary experiment. The geometric signature distinguishes training regimes but does not demonstrate predictive utility for robustness assessment without group labels, nor has it been validated on other spurious correlation benchmarks.
